\documentstyle[%


righttag%


,11pt%


,amssymb%


,russian%               remove in Eng.


,izvru%                 Set izven% in Eng.


]{amsart}





% 





\newcommand{\tts}[1]{}


\renewcommand{\baselinestretch}{0.98}





%\hoffset=-15mm%                  For Eng.


\setcounter{page}{3}


\god{2005}


\nomer{8~(519)} %                        Remove in Eng.


%\nomer{Vol.\,49, No.\,8}%               Set in Eng.


%\str{\pageref{begin}--\pageref{end}}%  Set in Eng.


\udk{517.521}





\author{\it А.В.\,Антонова}


% NB:   ^^ remove in Eng.


\title{Дополнение к признаку Жамэ}





\date{}


%\pagestyle{myheadings}%         Set in Eng.


\pagestyle{plain} %              Remove in Eng.


\begin{document}


%\thispagestyle{plain}%          Set in Eng.


\maketitle


\label{begin}








Целью данной работы является получение новых признаков сходимости  для


рядов с положительными членами,  которые позволяют 


упростить исследование таких рядов.





Пусть на множестве натуральных чисел ${\bold N}$ заданы три функции


$\varphi (n)$, $g(n)$, $f(n)$ такие, что


\begin{itemize}


\item[1)] $\varphi (n)$, $f(n)$ являются возрастающими, положительно


определенными, неограниченными функциями;


\item[2)] $g(n)$ положительно определена;


\item[3)] 


$\lim\limits_{n \to \infty} \frac{\varphi (n)g(n)}{f(n)\ln n}=q$;


\item[4)] $\lim\limits_{n \to \infty}\frac{g(n)}{f(n)}=0$.


\end{itemize}





Рассмотрим отдельно случаи, когда $q \ne 0$ --- конечное число, $q=0$


и $q =+\infty $.





Пусть $q \ne 0$ --- конечное число. Тогда при выше перечисленных


условиях справедлива








\begin{thm}


Если для элементов ряда $\sum\limits_{n=1}^{\infty}a_n$, начиная с


некоторого номера $N$, $\forall n \geq N$ выполняется


неравенство





$1)$ $(1-a_n^{ 1/\varphi(n)})f(n)/g(n)\geq p >1/q$,


где $p$ --- некоторое постоянное число,


то ряд сходится.





Если же $(\forall n \geq N)$ выполняется неравенство





$2)$ $(1-a_n^{1/\varphi(n)})f(n)/g(n)\leq p <1/q$,


где $p$ --- некоторое постояннное число, то ряд расходится.


\end{thm}








\begin{pf}


Пусть, начиная с некоторого номера


$N$, для любого $ n \geq N$ выполняется


неравенство 1) теоремы.


Тогда


\begin{equation}


(1-p g(n)/f(n)) ^{ \varphi (n)}\geq a_n.


\end{equation}


Докажем, что сходится ряд


\begin{equation}


\sum_{n=k}^{\infty}


(1-p g(n)/f(n))^{\varphi (n)},


\end{equation}


где $k$ --- номер, начиная с которого \medskip имеют смысл элементы ряда.


Применим к этому ряду логарифмический признак [1]


\begin{equation}


\frac{\ln(1-p g(n)/f(n))^{-\varphi(n)}}{\ln n}


=\frac{\varphi(n) p g(n)}{f(n) \ln n} \ln B(n),\quad


B(n)=(1-p g(n)/f(n))^{\frac{f(n)}{-p g(n)}}.


\end{equation}


В силу условия 4) получаем


$\lim\limits_{n \to \infty}\ln B(n)=1$.


Следовательно, для любого $\varepsilon >0$ найдется $N_1$ такое, что


для всех $n \geq N_1$ выполняется неравенство


\begin{equation}


1-\varepsilon <\ln B(n)< 1+\varepsilon.


\end{equation}


Из предположения 3) следует, что


для любого $\varepsilon ' >0$ существует $N_2$ такое, что для всех


$n \geq N_2$ выполняется неравенство


$$


q-\varepsilon ' < \frac{\varphi (n)g(n)}{f(n)\ln n}< q+\varepsilon '.


$$


Тогда получаем


$$


p \frac{\varphi (n) g(n)}{f(n)\ln n} \ln B(n)\geq


p (q-\varepsilon q - \varepsilon '+\varepsilon ' \varepsilon).


$$


Так как $p\geq \frac{1}{q-\delta}$,


где $0<\delta <q$ --- некоторое число, а числа $\varepsilon$ и


$\varepsilon '$ могут быть любыми, то выбираем их


таким образом, чтобы выполнялось неравенство


$q-\varepsilon q -\varepsilon '+


\varepsilon ' \varepsilon \geq q-\beta$,


где $0< \beta < \delta $ --- любое число.


Тогда будем иметь неравенство


$$


p(q-\varepsilon q -\varepsilon '+


\varepsilon ' \varepsilon) \geq \frac{q-\beta}{q-\delta}=\alpha >1,


$$


которое выполняется для всех $n \geq \max(N, N_1, N_2)$. В силу


логарифмического признака [1] ряд (2)


сходится, а значит, по признаку


сравнения [2] сходится и ряд $\sum\limits_{n=1}^{\infty}a_n$.





Пусть теперь для всех $n \geq N$ выполняется неравенство 2) теоремы.


Тогда


$$


(1-p g(n)/f(n)) ^{\varphi (n)}\leq a_n.


$$





Докажем, что расходится ряд (2). Применяя к этому ряду


логарифмический признак и рассуждая так же, как и при доказательстве


первой части, получим


$$


\frac{\ln(1-p g(n)/f(n))^{-\varphi (n)}}{\ln n}


=\frac{\varphi (n)p g(n)}{f(n) \ln n}


\ln B(n)\leq p(q+\varepsilon q+\varepsilon '+\varepsilon \varepsilon ').


$$





Так как $p\leq \frac{1}{q+\gamma}$,


где $\gamma >0$ --- некоторое число,


то, выбрав числа $\varepsilon $ и $\varepsilon '$ так, что


$q+\varepsilon q+


\varepsilon '+\varepsilon ' \varepsilon \leq q+\beta$,


где $0<\beta <\gamma$ --- любое число (такой выбор возможен 


в силу произвольности чисел $\varepsilon $ и


$\varepsilon '$), получим неравенство


$p (q+\varepsilon q +\varepsilon ' +\varepsilon ' \varepsilon ) \leq


\frac{q+\beta}{q+\gamma}=\alpha <1$,


которое выполняется для всех $n \geq M$, где $M\geq N$ --- некоторый


номер. В силу логарифмического признака ряд (2) расходится,


а значит, расходится и ряд


$\sum\limits_{n=1}^{\infty}a_n$.


\end{pf}





Рассмотрим теперь случай, когда $q=+\infty $. Для него имеет место








\begin{thm}


Если для элементов ряда


$\sum\limits_{n=1}^{\infty}a_n$, начиная с некоторого номера


$N$, $\forall n \geq N$ выполняется неравенство





$1)$ $(1-a_n^{1/\varphi (n)})f(n)/g(n)\geq p>0$, 


где $p$ --- некоторое постоянное число, то ряд сходится.





Если же $(\forall n \geq N)$ выполняется неравенство





$2)$ $(1-a_n^{1/ \varphi (n)})


f(n)/g(n)<0, $ то ряд расходится.


\end{thm}








\begin{pf}


Пусть, начиная с некоторого номера


$N$, $\forall n \geq N$ выполняется неравенство 1) теоремы. Тогда


имеет место ~(1).


Докажем, что сходится ряд (2).


Применяя к этому ряду логарифмический признак, получим равенство (3).





В силу условия 4) получаем, что


для любого $\varepsilon >0 $ найдется $N_1$ такое, что для всех


$n \geq N_1$ выполняется неравенство (4).


По условию~3) имеем


$\frac{\varphi (n) g(n)}{f(n) \ln n} \to +\infty $ при


$n \to \infty$,


следовательно, для каждого $A>0$ найдется $N_2$ такое, что для всех


$n \geq N_2$


$$


\frac{\varphi (n)g(n)}{f(n) \ln n}\geq A.


$$


Отсюда получаем


$p\frac{\varphi (n) g(n)}{f(n) \ln n} \ln B(n)\geq p A (1-\varepsilon)$.


Так как числа $A$ и $\varepsilon$ произвольны,


то выберем их таким образом, что


$A (1-\varepsilon) > \frac{1}{p-\delta}$,


где $0<\delta <p$ --- любое число.


Тогда будем иметь неравенство


$p A (1-\varepsilon )\geq \frac{p}{p-\delta}=\alpha >1$ для всех


$n \geq \max (N, N_1, N_2)$.


Следовательно, ряд (2)


сходится, а значит, по признаку сравнения сходится и ряд


$\sum\limits_{n=1}^{\infty} a_n$.





Пусть теперь для всех $ n \geq N$ выполняется неравенство 2) теоремы.


Из условий 1) и 2) получаем


$f(n)/g(n)>0$ для всех $n \geq N$. Следовательно,


$1-a_n^{1/\varphi(n)}<0$ \; $\forall n \geq N$, а значит,


$a_n>1$, которое влечет расходимость ряда


$\sum\limits_{n=1}^{\infty}a_n$.


\end{pf}





Наконец, в случае, когда $q=0$, имеет место








\begin{thm}


Если для элементов ряда 


$\sum\limits_{n=1}^{\infty}a_n$, начиная с некоторого номера


$N$, $\forall n \geq N$ выполняется неравенство


\begin{equation}


( 1-a_n^{1/\varphi (n)})f(n)/g(n)\leq p,


\end{equation}


где $p$ --- некоторое число,


то ряд расходится.


\end{thm}








\begin{pf}


Пусть, начиная с некоторого номера $N$,


для всех $n \geq N$ выполнено неравенство (5).


Тогда


$(1-p g(n)/f(n))^{\varphi (n)}\leq a_n$.


Докажем, что расходится ряд (2).


Применяя к нему логарифмический признак, получим равенство (3).





Поскольку последовательность


$\{\frac{\varphi (n) g(n)}{f(n) \ln n}\} \to 0$


при $n \to \infty$, а последовательность


$\{p\ln B(n)\}$ ограниченная, то последовательность


$\{ p\frac{\varphi (n) g(n)}{f(n) \ln n} \ln B(n)\}\to 0$


при $n \to \infty$. Отсюда следует, что расходится ряд (2),


а значит, по признаку сравнения расходится


и ряд $\sum\limits_{n=1}^{\infty}a_n$.


\end{pf}





{\bf Пример.} Исследовать на сходимость ряд


$$


\sum_{n=1}^{\infty}\bigg( \frac{a_k n^k+a_{k-1} n^{k-1}+\dotsb +a_0}{


b_kn^k +b_{k-1}n^{k-1}+\dotsb +b_0} \bigg)


^{c_sn^s+c_{s-1}n^{s-1}+\dotsb +c_0},


$$


где $a_i,b_i\geq 0$, $i=\overline{1,k}$, $c_j\geq0$,


$j=\overline{1,s}$.





Рассмотрим случай, когда $s>k>0$. Тогда, полагая


\begin{align*}


\varphi (n)&=c_sn^s+c_{s-1}n^{s-1}+\dotsb +c_0,\\


f(n)&=b_kn^k +b_{k-1}n^{k-1}+\dotsb +b_0,\\


g(n)&=1


\end{align*}


и применяя теорему 2 в предельной форме, получаем, что


\begin{itemize}


\item[1)] при $b_k>a_k$ ряд сходится;


\item[2)] при $b_k<a_k$ ряд расходится;


\item[3)] при $b_k=a_k$ и $b_{k-1}>a_{k-1}$ ряд сходится;


\item[\phantom{3)}] при $b_k=a_k$ и $b_{k-1}<a_{k-1}$ ряд расходится;


\item[\phantom{3)}] при $b_k=a_k$, $b_{k-1}=a_{k-1}$ и $b_{k-2}>a_{k-2}$


ряд сходится;


\item[\phantom{3)}] при $b_k=a_k$, $b_{k-1}=a_{k-1}$ и $b_{k-2}<a_{k-2}$


ряд расходится;


\item[\phantom{3)}] $\dotsc\dotsc\dotsc\dotsc\dotsc\dotsc\dotsc$


\item[\phantom{3)}] при $b_0>a_0$ ряд сходится,


\item[\phantom{3)}] при $b_0 \leq a_0$ ряд расходится.


\end{itemize}


Аналогично рассматриваются случаи, когда $s=k>0$ и $0<s<k$.








\begin{thebibliography}{9}





\bibitem{1}


Шмелев П.А. {\em Теория рядов в задачах и упражнениях}. Уч. пособие


для втузов. -- М.: Высш. школа, 1983. -- 176c.


\bibitem{2}


Фихтенгольц Г.М. {\em Курс дифференциального и интегрального исчисления}.


Т.\,2. -- М.: Наука, 1969. -- 800~c.





\end{thebibliography}





\vskip 0.5cm





\post{\it Казанский государственный}{\it Поступила}


\post{\it энергетический университет}{04.04.2003}





\label{end}


\end{document}


